\section{Related Work}
% 传统的LLM量化研究主要围绕INT8/INT4等整数格式展开，其核心在于通过激活感知策略来减小量化精度损失~\cite{illm,mambaquant}。例如，GPTQ通过量化感知微调来调整权重分布，以减轻异常值的影响~\cite{gptq,gptaq}。AWQ提出了一种激活感知的权重量化策略，基于激活统计来保护重要的权重通道~\cite{awq}。SmoothQuant则通过引入可学习的缩放因子，在激活和权重之间重新分配动态范围，从而缓解整数量化中的溢出问题~\cite{smoothquant}。然而，这些方法本质上是在对抗“固定范围”的限制，未能利用浮点格式的动态指数特性。
Traditional research on LLM quantization has primarily revolved around integer formats such as INT8 and INT4, focusing on reducing quantization accuracy loss through activation-aware strategies~\cite{illm,mambaquant}. For instance, GPTQ adjusts weight distributions via quantization-aware updates to mitigate the impact of outliers~\cite{gptq,gptaq}. AWQ proposes an activation-aware weight quantization strategy that protects salient weight channels based on activation statistics~\cite{awq}. SmoothQuant introduces scaling factors to redistribute the dynamic range between activations and weights, thereby alleviating overflow issues inherent to integer quantization~\cite{smoothquant}. However, these methods fundamentally combat the limitations of "fixed ranges" and fail to exploit the dynamic exponent characteristics of floating-point formats.

% 近年来，旋转技术成为量化领域的热点，并成为当前W4A4量化的SOTA方向。QuIP#~\cite{quip_sharp}和QuaRot~\cite{quarot} 引入了 Hadamard 变换以平滑权重和激活的分布；OSTQuant和 FlatQuant虽然在 W4A4 整数量化上取得了 SOTA，但其目标函数旨在压缩数据的动态范围（Clip Ratio Optimization），这与 MXFP4 需要充分利用动态范围的特性并不相符。此外，BRQ~\cite{brq} 虽然尝试了针对 MXFP4 的旋转，但其仅关注块内分布的简单调整，忽略了块间能量失衡这一关键因素。
Recently, rotation techniques have emerged as a focal point in quantization research, representing the current state-of-the-art (SOTA) direction for W4A4 quantization. QuIP\#~\cite{quip_sharp} and QuaRot~\cite{quarot} introduce the Hadamard transform to smooth the distributions of weights and activations. While OSTQuant and FlatQuant have achieved SOTA results in W4A4 integer quantization, their objective functions aim to compress data dynamic range (Clip Ratio Optimization). This objective is incompatible with MXFP4, which requires fully utilizing dynamic range. Furthermore, although BRQ~\cite{brq} attempts rotation for MXFP4, it focuses solely on simple adjustments of intra-block distribution, neglecting the critical factor of inter-block energy imbalance.

% 本文提出的TORQ框架，是首个针对块浮点（Block Floating Point）结构特征设计的双级旋转框架。本文通过“宏观均衡旋转”与“微观对齐旋转”相结合系统性解决了MXFP4激活量化的核心挑战。这一工作填补了针对块浮点结构进行跨块与块内联合分布优化的空白，为低比特浮点格式的激活量化提供了新的优化范式。
The TORQ framework proposed in this paper is the first dual-level rotation framework explicitly designed for the structural characteristics of Block Floating Point (BFP). By combining "Macro-Equilibrium Rotation" and "Micro-Alignment Rotation," we systematically address the core challenges of MXFP4 activation quantization. This work addresses the lack of in joint inter-block and intra-block distribution optimization for block floating-point structures, establishing a new optimization paradigm for low-bit floating-point activation quantization.